\frfilename{mt28.tex} 
\versiondate{17.1.15} 
\copyrightdate{1994} 
 
\def\chaptername{Analisis Fourier} 
\def\sectionname{Pendahuluan} 
 
\newchapter{28} 
 
Untuk bab terakhir jilid ini, saya mencoba menyajikan uraian singkat tentang salah satu topik terpenting dalam analisis. Ini adalah usaha yang berani, dan saya tidak dapat berharap memenuhi tuntutan wajar siapa pun yang mengenal dan mencintai subjek ini sebagaimana mestinya. Namun saya juga tidak dapat melewatinya tanpa berlaku tidak setia pada bidang saya sendiri, karena persoalan-persoalan yang disumbangkan oleh kajian deret dan transformasi Fourier telah mengarahkan teori ukuran sepanjang sejarahnya. Oleh karena itu, yang akan saya coba lakukan adalah memberikan versi hasil-hasil tersebut yang patut diketahui semua orang, dalam bahasa yang menyatukannya dengan bagian-bagian lain dari risalah ini, dengan tujuan membuka saluran perpindahan intuisi dan teknik antara kajian abstrak umum tentang ruang ukur, yang menjadi pusat karya kita, dan keluarga penerapan teori integrasi yang khusus ini. 
 
Saya membagi materi bab ini, secara cukup konvensional, menjadi tiga bagian: deret Fourier, transformasi Fourier, dan fungsi karakteristik dalam teori probabilitas. Walaupun jelas bahwa banyak gagasan yang sama berlaku bagi ketiganya, pada tahap ini saya tidak menganggap berguna untuk mencoba merumuskan generalisasi eksplisit yang menyatukan ketiganya; hal itu termasuk teori analisis harmonik yang lebih umum pada grup, yang harus menunggu sampai Jilid 4. Namun saya memulai dengan seksi tentang teorema Stone-Weierstrass (\S281), yang merupakan salah satu alat dasar analisis fungsional, sekaligus berguna untuk bab ini. Seksi terakhir (\S286), yakni bukti teorema Carleson, berada pada tingkat yang agak berbeda dari bagian-bagian lainnya. 
 
\discrpage 
 
